% TeachingMaterials/Layouts/Scripts/inc/commands.tex

%% ============================================================================ %%
%% Custom commands
%% ============================================================================ %%

\newcommand{\cb}{\colorbox}
\newcommand{\tc}{\textcolor}
\newcommand{\tb}{\textbf}
\newcommand{\p}{\partial}
\newcommand{\te}{\text}
\newcommand{\f}{\frac}
%% ============================================================================ %%


%% ============================================================================ %%
%% Theorem environment
%% ============================================================================ %%

\mdfdefinestyle{theoremstyle}{%
  backgroundcolor=dblue!2,                          % background color
  linecolor=dblue,                                  % Border color
  linewidth=2pt,                                    % Border thickness
  rightline=true,                                   % Right border visible
  leftline=true,                                    % Left border visible
  topline=false,                                    % Top border visible
  bottomline=false,                                 % Bottom border visible
  innertopmargin=10pt,                              % Margin inside the top
  innerbottommargin=10pt,                           % Margin inside the bottom
  innerleftmargin=10pt,                             % Margin inside the left
  innerrightmargin=10pt,                            % Margin inside the right
  frametitlerule=true,                              % Title rule visible
  frametitlerulewidth=0pt,                          % Thickness of title rule
  frametitlerulecolor=dblue,                        % Color of title rule
  frametitlebackgroundcolor=dblue!5,                % Background color for the title
  roundcorner=10pt                                  % Rounded corners
}

% == Define a theorem counter == %
\newtheorem{innerthm}{Theorem}[chapter]

% == Define a custom theorem environment == %
\newenvironment{thm}[1][]{%
  \begin{mdframed}[style=theoremstyle,frametitle={\textcolor{dblue}{Theorem~\theinnerthm\ifx\\#1\\\else: #1\fi}}]%
  \refstepcounter{innerthm}                         % Increment the theorem counter
  \vspace{-0.5em}                                   % Adjust spacing
}{%
  \end{mdframed}%
}
%% ============================================================================ %%


%% ============================================================================ %%
%% Definition environment
%% ============================================================================ %%

\mdfdefinestyle{definitionstyle}{%
  backgroundcolor=porange!2,                        % background color
  linecolor=porange!95!black,                                % Border color
  linewidth=2pt,                                    % Border thickness
  rightline=true,                                   % Right border visible
  leftline=true,                                    % Left border visible
  topline=false,                                    % Top border visible
  bottomline=false,                                 % Bottom border visible
  innertopmargin=10pt,                              % Margin inside the top
  innerbottommargin=10pt,                           % Margin inside the bottom
  innerleftmargin=10pt,                             % Margin inside the left
  innerrightmargin=10pt,                            % Margin inside the right
  frametitlerule=true,                              % Title rule visible
  frametitlerulewidth=0pt,                          % Thickness of title rule
  frametitlerulecolor=porange,                      % Color of title rule
  frametitlebackgroundcolor=porange!5,              % Background color for the title
  roundcorner=10pt                                  % Rounded corners
}

% == Define a theorem counter == %
\newtheorem{innerdef}{Definition}[chapter]

% Define a custom theorem environment
\newenvironment{definition}[1][]{%
\refstepcounter{innerdef} % Increment the Definition counter
\begin{mdframed}[style=definitionstyle,frametitle={\textcolor{porange!95!black}{Definition~\theinnerdef\ifx\\#1\\\else: #1\fi}}]%
\vspace{-0.5em} % Adjust spacing
}{%
\end{mdframed}%
}
%% ============================================================================ %%


%% ============================================================================ %%
%% Proof environment
%% ============================================================================ %%

\mdfdefinestyle{proofstyle}{%
backgroundcolor=dblue!2,                          % background color
linecolor=dblue!30!lightgray,                     % Border color
linewidth=2pt,                                    % Border thickness
rightline=true,                                   % Right border visible
leftline=true,                                    % Left border visible
topline=false,                                    % Top border visible
bottomline=false,                                 % Bottom border visible
innertopmargin=10pt,                              % Margin inside the top
innerbottommargin=10pt,                           % Margin inside the bottom
innerleftmargin=10pt,                             % Margin inside the left
innerrightmargin=10pt,                            % Margin inside the right
frametitlerule=true,                              % Title rule visible
frametitlerulewidth=0pt,                          % Thickness of title rule
frametitlerulecolor=dblue!30!lightgray,           % Color of title rule
frametitlebackgroundcolor=dblue!5,                % Background color for the title
roundcorner=10pt                                  % Rounded corners
}

% == Define a theorem counter == %
\newtheorem{innerproof}{Proof}[chapter]

% Define a custom theorem environment
\newenvironment{proof}[1][]{%
  \begin{mdframed}[style=proofstyle,frametitle={\textcolor{dblue!30!lightgray}{Proof of Theorem~\ref{#1}}}]%
  \vspace{-0.5em} % Adjust spacing
}{%
  \end{mdframed}%
}
%% ============================================================================ %%


%% ============================================================================ %%
%% Example environment
%% ============================================================================ %%

\mdfdefinestyle{examplestyle}{%
backgroundcolor=lightgray!2,                      % background color
linecolor=lightgray,                              % Border color
linewidth=2pt,                                    % Border thickness
rightline=true,                                   % Right border visible
leftline=true,                                    % Left border visible
topline=false,                                    % Top border visible
bottomline=false,                                 % Bottom border visible
innertopmargin=10pt,                              % Margin inside the top
innerbottommargin=10pt,                           % Margin inside the bottom
innerleftmargin=10pt,                             % Margin inside the left
innerrightmargin=10pt,                            % Margin inside the right
frametitlerule=true,                              % Title rule visible
frametitlerulewidth=0pt,                          % Thickness of title rule
frametitlerulecolor=lightgray,                    % Color of title rule
frametitlebackgroundcolor=lightgray!5,            % Background color for the title
roundcorner=10pt                                  % Rounded corners
}

% == Define a theorem counter == %
\newtheorem{innerexpl}{Example}[chapter]

% Define a custom theorem environment
\newenvironment{example}[1][]{%
\refstepcounter{innerexpl} % Increment the Definition counter
\begin{mdframed}[style=examplestyle,frametitle={\textcolor{lightgray!95!black}{Example~\theinnerdef\ifx\\#1\\\else: #1\fi}}]%
\vspace{-0.5em} % Adjust spacing
}{%
\end{mdframed}%
}
%% ============================================================================ %%


%% ============================================================================ %%
%% Code environment
%% ============================================================================ %%

\mdfdefinestyle{codestyle}{%
backgroundcolor=trf!2,                            % background color
linecolor=trf!70!,                                % Border color
linewidth=2pt,                                    % Border thickness
rightline=true,                                   % Right border visible
leftline=true,                                    % Left border visible
topline=false,                                    % Top border visible
bottomline=false,                                 % Bottom border visible
innertopmargin=10pt,                              % Margin inside the top
innerbottommargin=10pt,                           % Margin inside the bottom
innerleftmargin=10pt,                             % Margin inside the left
innerrightmargin=10pt,                            % Margin inside the right
frametitlerule=true,                              % Title rule visible
frametitlerulewidth=0pt,                          % Thickness of title rule
frametitlerulecolor=trf!70!,                      % Color of title rule
frametitlebackgroundcolor=trf!5,                  % Background color for the title
roundcorner=10pt,                                 % Rounded corners
font=\ttfamily                                    % Set font to monospace
}

% == Define a theorem counter == %
\newtheorem{innercode}{Code}[chapter]

% Define a custom theorem environment
\newenvironment{code}[1][]{%
\refstepcounter{innercode} % Increment the Definition counter
\begin{mdframed}[style=codestyle,frametitle={\textcolor{trf!70!}{Code~\theinnercode\ifx\\#1\\\else: #1\fi}}]%
\vspace{-0.5em} % Adjust spacing
}{%
\end{mdframed}%
}
%% ============================================================================ %%



%% ============================================================================ %%
%% Important environment
%% ============================================================================ %%

\mdfdefinestyle{impstyle}{%
  backgroundcolor=porange!2,                        % background color
  linecolor=porange!95!black,                       % Border color
  linewidth=2pt,                                    % Border thickness
  rightline=true,                                   % Right border visible
  leftline=true,                                    % Left border visible
  topline=false,                                    % Top border visible
  bottomline=false,                                 % Bottom border visible
  innertopmargin=10pt,                              % Margin inside the top
  innerbottommargin=10pt,                           % Margin inside the bottom
  innerleftmargin=10pt,                             % Margin inside the left
  innerrightmargin=10pt,                            % Margin inside the right
  frametitlerule=true,                              % Title rule visible
  frametitlerulewidth=0pt,                          % Thickness of title rule
  frametitlerulecolor=porange,                      % Color of title rule
  frametitlebackgroundcolor=porange!5,              % Background color for the title
  roundcorner=10pt                                  % Rounded corners
}

% == Define a theorem counter == %
%\newtheorem{innerdef}{Definition}[chapter]

% Define a custom theorem environment
\newenvironment{imp}[1][]{%
%\refstepcounter{innerdef} % Increment the Definition counter
\begin{mdframed}[style=impstyle,frametitle={\textcolor{porange!95!black}{\ifx\\#1\\\else #1\fi}}]%
\vspace{-0.5em} % Adjust spacing
}{%
\end{mdframed}%
}
%% ============================================================================ %%


%% ============================================================================ %%
%% Custom enumerate environment
%% ============================================================================ %%

% Customize enumerate with vertical bar
%\setlist[enumerate,1]{label=\textcolor{lightgray}{\arabic*} \textcolor{lightgray}{|}, labelsep=0.5em, left=0em}
\setlist[enumerate,1]{%
  label=\textcolor{lightgray}{\arabic*} \hspace{0em} {\textcolor{lightgray}{\rule[-0.2ex]{1pt}{2ex}}}, % Number + vertical bar as a rule
  labelsep=0.75em, % Space between label (number + bar) and text
  left=0em % Remove additional indentation
}

\setlist[enumerate,2]{%
  label=\textcolor{lightgray}{\roman*} \hspace{0em} {\textcolor{lightgray}{\rule[-0.2ex]{1pt}{2ex}}}, % Number + vertical bar as a rule
  labelsep=0.75em, % Space between label (number + bar) and text
  left=0em % Remove additional indentation
}
%% ============================================================================ %%

